\section{О себе}
\subsection{Навыки/Опыт/известные алгоритмы}
Так как довольно просто выполнить тестовое на хорошем уровне из-за существования LLM, я вынужден написать подробный блок о себе

В данный момент я обучаюсь на потоке М+П+ ФКН ПМИ.

В силу отличной учебы на пилотном курсе АиСД-1 (Итоговая оценка: 9), оценка по АиСД-2 пока неизвестна, но я планирую так же получить отл. 9. я знаком с большинством стандартных продвинутых алгоритмов (в частности нужные для LSM дерева фильтр Блюма и Skip-List). Отдельно перечислю известные алгоритмы во внешней памяти: List ranking за $O(Sort(n))$, B/$B^+$-дерево, Куча во внешней памяти за $O(\frac1B \lceil \log_{M/B} \frac{q}{M} \rceil)$ на запрос, BFS во внешней памяти за $O(Sort(E) + V)$, поиск компонент связности за $O(Sort(E+V)\log V)$

Так же, я люблю много писать на С++ и достаточно хорошо им владею (закончил пилотный курс на 10, и продолжаю учится на продвинутом курсе сейчас). Есть опыт разработки больших проектов: написал свой интерепретатор для некоторого ЯП \cite{interpretator}. Так же, приведу ниже менее масштабные pet-проекты:

\begin{enumerate}
    \item Интерактивные деревья поиска \cite{bstgarden}
    \item Программа для обработки фото \cite{imageprocessor} (вообще говоря, это был общекурсовой проект, но я тут заморочился с ООП, так что считаю необходимым это здесь прикрепить), p.s. все файлы сгружены в кучу, потому что тест система иначе не принимала
\end{enumerate}


\subsection{Чем вызван интерес}

Мне очень понравились задачи во внешней памяти, которые нам предложили на курсе алгоритмов. И вообще меня заинтересовала разработка каких-то структур во внешней памяти, а данная тема для курсовой позволит мне реализовать данное желание
